\section{A Running Example}
\label{sec:example}
Consider a routine that replaces an owned integer by its least nonnegative
integer square root rounded upward. The contract relates the initial value
$x$, the returned value, and the final cell contents. Its imported Rocq
proposition is
\begin{equation}
\label{eq:least-root}
\begin{split}
\mathsf{LeastRoot}(n,r)\equiv{}&0\le r\ \land\\[-2pt]
 &\forall k\in\mathbb Z.\;0\le k\Rightarrow
       (r\le k\Leftrightarrow n\le k^2).
\end{split}
\end{equation}
Listing~\ref{lst:root} shows the C body and abbreviated annotations of the
archived instance. We use $p\mapsto_{\mathrm{int}}v$ for typed cell ownership;
the artifact uses QCP's \texttt{data\_at} notation.

\begin{lstlisting}[caption={Least-root instance with an intermediate observation. The imported relation is Equation~\eqref{eq:least-root}.},label={lst:root}]
int root_cell(int *p)
/*@ With x
    Require 0 <= x && x <= 100 &&
            data_at(p, int, x)
    Ensure least_root(x, __return) &&
           data_at(p, int, __return) */
{
  int n = *p;
  int r = 0;
  /*@ Inv n == x && 0 <= r && r <= 10 &&
      (r == 0 || (r-1)*(r-1) < n) &&
      data_at(p, int, x) */
  while (r*r < n) { r = r + 1; }
  /*@ Assert least_root(n, r) &&
              data_at(p, int, x) */
  *p = r;
  return r;
}
\end{lstlisting}

\paragraph{A logical instance.}
The binding $\theta=\{x\mapsto13\}$ fixes the logical input and retains the
address $p$. The spatial precondition describes the cell at that symbolic
address. Loop execution produces $r=0,1,2,3,4$ at its five invariant boundaries;
then the body writes 4 and returns 4. The final observation is
\[
\mathsf{LeastRoot}(13,4)\land p\mapsto_{\mathrm{int}}4.
\]
Local scalar storage is handled by the observation frame described in
Section~\ref{sec:boundaries}. The imported relation remains a proposition in
the verification logic. Automatic execution establishes the entry and safety
phases, and proof completion checks the remaining functional and observation
obligations against their original generated statements.

\paragraph{Testing an intermediate specification.}
Replacing the assertion by $\mathsf{LeastRoot}(n,r-1)$ changes the observation
while leaving the C body and final contract intact. At that boundary the
program still has $n=13,r=4$, so the candidate asks for
$\mathsf{LeastRoot}(13,3)$. Instantiating its universal quantifier with $k=3$
exposes the contradiction $13\le9$. The failed observation does not replace the reached state or alter the
following write. Both variants therefore follow the same program computation.

\paragraph{Original assertions and executable reformulations.}
The original relation ranges over all nonnegative integers even after
$n$ and $r$ are fixed. An executable reformulation is
\begin{equation}
\label{eq:finite-root}
0\le r\land n\le r^2\land
   \bigl(r=0\lor(r-1)^2<n\bigr).
\end{equation}
The two formulas are equivalent over mathematical integers; the C contract
bounds inputs to control machine arithmetic. Constructing and justifying the
finite form is one route to a runtime oracle. \tespec{} instead retains the
original relation as a proposition and reuses its proof library. The checked
comparison appears in Section~\ref{sec:case-studies}.

\begin{figure*}[t]
\centering
\begin{tikzpicture}[x=1cm,y=1cm,>=Latex,
 font=\small,
 box/.style={draw=black!45,rounded corners=2pt,align=center,inner sep=5pt},
 state/.style={box,fill=blue!4,draw=blue!55!black,text width=3.75cm,minimum height=.9cm},
 goal/.style={box,fill=teal!5,draw=teal!60!black,text width=3.75cm,minimum height=1.0cm},
 flow/.style={->,line width=.8pt},
 tag/.style={font=\footnotesize,fill=white,inner sep=2pt}]
\node[box,text width=14.9cm,fill=black!3] at (5.2,0)
 {\textbf{Testing and verification share symbolic execution and proof logic}\\
 $\mathsf{LeastRoot}(n,r)\equiv 0\le r\land
 \forall k\in\mathbb Z.\;0\le k\Rightarrow(r\le k\Leftrightarrow n\le k^2)$};
\node[anchor=west,font=\footnotesize\bfseries,text=blue!55!black] at (-2.3,-.83)
 {SYMBOLIC EXECUTION \quad logical input $n=13$; address $p$ remains symbolic};
\node[state] (entry) at (0,-1.6)
 {\textbf{Initial domain}\\$p\mapsto_{\rm int}13$};
\node[state] (loop) at (5.2,-1.6)
 {\textbf{Symbolic loop states}\\$r=0,1,2,3,4$};
\node[state] (ret) at (10.4,-1.6)
 {\textbf{Symbolic write and return}\\$p\mapsto_{\rm int}4,\quad\mathit{ret}=4$};
\draw[flow] (entry)--node[tag]{execute}(loop);
\draw[flow] (loop)--node[tag]{execute}(ret);
\node[goal] (exists) at (0,-3.25)
 {\textbf{Entry obligation}\\$\exists z.\,D_\theta(z)$};
\node[goal] (mid) at (5.2,-3.25)
 {\textbf{Intermediate obligations}\\visited invariants;\\$\mathsf{LeastRoot}(13,4)$};
\node[goal] (post) at (10.4,-3.25)
 {\textbf{Return obligation}\\$\mathsf{LeastRoot}(13,4)\land$\\$p\mapsto_{\rm int}4$};
\draw[flow] (entry)--node[tag,right]{instantiate}(exists);
\draw[flow] (loop)--node[tag,right]{observe copy}(mid);
\draw[flow] (ret)--node[tag,right]{observe copy}(post);
\node[box,text width=14.9cm,fill=teal!5,draw=teal!60!black] (oracle) at (5.2,-4.95)
 {\textbf{Reuse the verification oracle for functional correctness}\\
 evaluate computable terms $\longrightarrow$ automatic reasoning
 $\longrightarrow$ checked proofs for residual properties};
\draw[flow] (exists.south)--(oracle.north -| exists.south);
\draw[flow] (mid.south)--(oracle.north);
\draw[flow] (post.south)--(oracle.north -| post.south);
\node[anchor=north,align=center,font=\footnotesize,text width=15.0cm] at (5.2,-5.62)
 {Changing the intermediate assertion to $\mathsf{LeastRoot}(13,3)$ changes its proof obligation.\\
 The program still writes and returns 4. Quantified properties retain their original logical form.};
\end{tikzpicture}

\caption{Instance checking in \tespec{}. Binding $x=13$ retains a symbolic
address $p$; execution updates its cell from 13 to 4.
A dashed copy produces an obligation without changing the execution.
Automatic decisions and kernel-checked completion populate the VC record
(V: VALID, I: INVALID, ?: UNKNOWN); the automatic route has its own
soundness interface (Section~\ref{sec:semantics}).}
\label{fig:overview}
\end{figure*}

